\documentclass[10pt,letterpaper]{article}

\usepackage[margin=1in]{geometry}
\usepackage{amsmath,amssymb,amsfonts}
\usepackage{booktabs}
\usepackage{multirow}
\usepackage{graphicx}
\usepackage{xcolor}
\usepackage{hyperref}
\usepackage{natbib}
\usepackage{tabularx}
\usepackage{array}

\newcolumntype{C}{>{\centering\arraybackslash}X}
\newcolumntype{R}{>{\raggedleft\arraybackslash}X}

\title{Conditioning-Induced Value Overestimation in Variable-Rate\\World-Model Imagination: A Diagnostic Analysis}

\author{
Anonymous Authors
}

\date{}

\begin{document}

\maketitle

%% ============================================================
%% ABSTRACT
%% ============================================================
\begin{abstract}
Temporal abstraction---imagining at variable time scales---is a natural capability of learned world models, yet prior attempts to exploit it in DreamerV3 consistently degrade performance.
We present a systematic factorial diagnosis identifying the root cause: a \emph{trained dt-embedding} that conditions the critic on temporal-skip information triggers \textbf{Conditioning-Induced Value Overestimation (CIVO)}, a positive-feedback loop in which the critic exploits precise but inaccurate dt-conditioned predictions from the RSSM.
Using a normalized value-to-eval (NVE) diagnostic, we show that a 1.5$\times$ threshold at 35K+ steps perfectly separates pathological from healthy runs (16/16 accuracy across 4 environments: Walker Walk, Cheetah Run, Cartpole Swingup, Hopper Hop).
Gradient decomposition reveals that 93\% of the dt-embedding gradient originates from observation decoders, not the reward head, pinpointing the architectural vulnerability.
Restricting imagination to dt$=$1 eliminates the primary CIVO pathway but exposes a secondary eval-driven instability (Pathway B, 1/4 seeds), demonstrating that surface-level mitigations are insufficient.
Notably, TAWM (Nhu et al., 2025) succeeds with dt-conditioning under MPC, confirming that the vulnerability is specific to imagination-based actor-critic architectures.
Our findings provide actionable guidance: do not train dt-embeddings for temporal conditioning in imagination-based actor-critic systems without safeguards against the observation-decoder gradient pathway.
\end{abstract}

%% ============================================================
%% 1. INTRODUCTION
%% ============================================================
\section{Introduction}

Temporal abstraction---the ability to reason and plan at multiple time scales---has long been a goal of reinforcement learning research, from the options framework \citep{sutton1999options} to macro-actions and action repetition \citep{lakshminarayanan2017dynamic}.
World-model reinforcement learning, exemplified by the Dreamer family \citep{hafner2020dreamer,hafner2023dreamerv3}, offers a particularly natural testbed: one can modify the imagination horizon's temporal granularity without changing environment interaction, enabling variable-rate ``mental simulation.''

Despite this appeal, our systematic exploration of variable-rate imagination in DreamerV3 reveals a striking negative result: \emph{every} configuration that introduces dt$>$1 imagination degrades asymptotic performance relative to the dt$=$1 baseline.
Fixed dt$=$2 imagination peaks at 485 (51\% of baseline 950) before collapsing; dt$=$4 reaches only 177 (19\% of baseline); learned gating peaks at 541 then collapses to 99.

This paper asks: \emph{why does temporal abstraction fail in world-model RL?}

Through a controlled factorial experiment design, we identify the root cause as \textbf{Conditioning-Induced Value Overestimation (CIVO)}: when the dt-embedding is \emph{trained} (receives gradient from the world model's observation decoders), it enables the critic to make precise but systematically biased dt-specific value predictions.
The RSSM's prediction fidelity degrades at dt$>$1 (reward gap 2.78$\times$ for trained dt-embedding vs.\ 1.01$\times$ without), but the trained embedding creates the illusion of reliable conditioning, triggering a positive-feedback loop: overestimated values $\to$ actor exploitation $\to$ policy degradation $\to$ further overestimation.

Our contributions are:
\begin{enumerate}
    \item A \textbf{complete causal chain} from trained dt-embedding $\to$ observation-decoder gradient amplification $\to$ CIVO $\to$ performance collapse, validated across 4 DMC environments with 16/16 NVE threshold accuracy.
    \item A \textbf{gradient decomposition} showing 93\% of dt-embedding gradient comes from observation decoders (not reward), identifying the precise architectural vulnerability and explaining why reward-head interventions are ineffective.
    \item A \textbf{dual-pathway characterization}: CIVO (Pathway A, value-driven, entropy$\downarrow$) vs.\ eval-driven instability (Pathway B, entropy$\uparrow$), showing that trained dt-embedding creates vulnerability through multiple mechanisms.
    \item A \textbf{cross-architecture contrast}: dt-conditioning succeeds in MPC (TAWM) but fails in imagination-based actor-critic, localizing the vulnerability to the imagination training loop.
\end{enumerate}

%% ============================================================
%% 2. BACKGROUND AND RELATED WORK
%% ============================================================
\section{Background and Related Work}

\paragraph{Temporal abstraction in RL.}
The options framework \citep{sutton1999options} formalizes temporal abstraction through temporally extended actions.
FiGAR \citep{sharma2017figar} learns action repetition counts.
Action repeat is widely used as a hyperparameter in continuous control \citep{hafner2023dreamerv3}.
These model-free approaches do not suffer from the world-model fidelity gap that drives CIVO.

\paragraph{World-model RL.}
The Dreamer family \citep{hafner2020dreamer,hafner2021dreamerv2,hafner2023dreamerv3} learns latent-dynamics models (RSSM) and trains actor-critic networks entirely in imagination.
IRIS \citep{micheli2023iris} and TWM \citep{robine2023twm} use transformer-based world models.
Variable-rate imagination has been explored informally but without systematic diagnosis of failure modes.

\paragraph{Temporal abstraction in world models.}
TAWM \citep{nhu2025tawm} successfully uses dt-conditioning in TD-MPC2, an MPC-based planner.
Critically, TAWM does \emph{not} train actor-critic networks in imagination---it uses the world model for online planning.
Our finding that dt-conditioning fails specifically in imagination-based actor-critic training but succeeds in MPC is a key novelty.

\paragraph{Value overestimation in MBRL.}
Model exploitation and value overestimation are known challenges in model-based RL \citep{janner2019mbpo}.
Lu et al.\ (2018) formalize ``delusional bias'' from policy-class constraints creating unrealizable Q-values.
CIVO is mechanistically distinct: it arises from \emph{model imprecision amplified by learned conditioning}, not from policy-class constraints.
We use the term CIVO to avoid confusion with Lu et al.'s delusional bias.

%% ============================================================
%% 3. METHODOLOGY
%% ============================================================
\section{Methodology}

\subsection{Base Architecture}

We use DreamerV3 \citep{hafner2023dreamerv3} with its default RSSM world model, actor-critic imagination training, and symlog predictions.
The environment interaction loop is unmodified; all changes occur within the imagination and world-model training pipelines.

\subsection{Variable-Rate Imagination}

We introduce a \textbf{dt-embedding} $e_{\text{dt}} \in \mathbb{R}^d$ that conditions the RSSM transition model, reward head, and continuation head on the temporal skip $\text{dt} \in \{1,2,4\}$.
\textbf{Temporal subsampling} constructs training sequences by skipping every other transition (dt$=$2) or every third (dt$=$4), with the RSSM learning to predict across the gap.
A \textbf{learned gating network} $g_\phi(s_t) \in [0,1]$ selects dt per imagination step.

\subsection{Factorial Experiment Design}

We design a controlled ablation isolating the factors contributing to failure:

\begin{itemize}
    \item \textbf{dt-embedding status}: trained (gradient ON, subsample ON), untrained (gradient ON, subsample OFF---never sees dt$>$1 data), frozen (gradient OFF, random init), absent (dim$=$0).
    \item \textbf{Temporal subsampling}: ON (dt$>$1 transitions in training data) vs.\ OFF.
    \item \textbf{Imagination horizon}: H$=$8 vs.\ H$=$15.
    \item \textbf{Imagination dt}: dt$=$1 only vs.\ dt$=$2 vs.\ learned gate.
\end{itemize}

\subsection{Diagnostic Metrics}

\paragraph{Normalized Value-to-Eval ratio (NVE).}
We define $\text{NVE}(t) = \frac{v(t)/e(t)}{v_{\text{base}}(t)/e_{\text{base}}(t)}$ where $v(t)$ is the mean predicted value and $e(t)$ is the evaluation return at step $t$.
NVE$>$1.0 indicates the value function is more optimistic than baseline calibration; NVE$>$1.5$\times$ at $t \geq$ 35K (when baseline v/e has converged past 0.24) serves as the pathological threshold.
The 35K+ constraint prevents false positives from early-training noise when v/e denominators are small.

\paragraph{Actor entropy.}
Policy entropy distinguishes CIVO (Pathway A: entropy decreases as actor becomes overconfidently deterministic) from eval-driven instability (Pathway B: entropy increases as policy destabilizes toward randomness).

\paragraph{dt-embedding cosine similarity.}
$\cos(e_{\text{dt}=1}, e_{\text{dt}=2})$ tracks embedding divergence during training, providing a mechanistic indicator of when the critic gains dt-discriminative power.

\subsection{Environments}

We evaluate across four DeepMind Control Suite environments spanning different observation dimensionalities and dynamics:
Walker Walk (24-dim), Cheetah Run (17-dim), Cartpole Swingup (5-dim), and Hopper Hop (15-dim).
Cross-environment replication elevates findings from task-specific to architectural.

%% ============================================================
%% 4. EXPERIMENTS AND RESULTS
%% ============================================================
\section{Experiments and Results}

\subsection{All Variable-Rate Configurations Degrade Performance (C1)}

Table~\ref{tab:c1} summarizes all dt$>$1 configurations tested in Walker Walk.
Every configuration degrades relative to the dt$=$1 baseline ($\sim$950 eval return), with no positive result for variable-rate imagination.

\begin{table}[h]
\centering
\caption{Variable-rate imagination consistently degrades performance in Walker Walk.}
\label{tab:c1}
\begin{tabular}{lccl}
\toprule
\textbf{Configuration} & \textbf{Peak Eval} & \textbf{\% of Baseline} & \textbf{Outcome} \\
\midrule
Baseline (dt$=$1) & 950 & 100\% & Stable convergence \\
Fixed dt$=$2 (H$=$15, trained) & 485 & 51\% & Collapse to $\sim$135 \\
Fixed dt$=$4 & 177 & 19\% & Severe degradation \\
Learned gate & 541 & 57\% & Collapse to 99 \\
\bottomrule
\end{tabular}
\end{table}

\subsection{Trained dt-Embedding Triggers Pathological Value Overestimation (C2)}
\label{sec:c2}

The critical ablation compares four dt-embedding configurations in Walker Walk (Table~\ref{tab:ablation}).
Only the trained dt-embedding produces NVE$>$1.5$\times$; all other configurations remain below.

\begin{table}[h]
\centering
\caption{2$\times$2 dt-embedding ablation matrix (Walker Walk, H$=$8). Only trained dt-embedding with active gradient flow produces pathological NVE.}
\label{tab:ablation}
\begin{tabular}{lcccc}
\toprule
\textbf{Config} & \textbf{dt\_emb Status} & \textbf{NVE Peak (35K+)} & \textbf{Collapse} & \textbf{Failure Mode} \\
\midrule
008d & Trained (grad ON) & 4.12$\times$ & Yes (1/2) & CIVO \\
008f & Absent (dim$=$0) & 1.26$\times$ & No & Reward degradation \\
015b & Frozen (grad OFF) & 0.47$\times$ & No & Gradient starvation \\
008c & Untrained (sub OFF) & 0.90$\times$ & No & Learning signal reduction \\
\bottomrule
\end{tabular}
\end{table}

\paragraph{Multi-seed NVE trajectories.}
Table~\ref{tab:nve_seeds} shows the NVE trajectory for the trained dt-embedding configuration (008d) across two seeds.
Both seeds exceed 1.5$\times$ NVE at 35K+, but outcomes diverge: seed 1 sustains overestimation and collapses (eval 131@105K), while seed 2 self-corrects (eval 268@105K).

\begin{table}[h]
\centering
\caption{NVE trajectory comparison for trained dt-embedding (008d), Walker Walk. Both seeds cross the 1.5$\times$ threshold; collapse is stochastic.}
\label{tab:nve_seeds}
\begin{tabular}{ccccc}
\toprule
\textbf{Step} & \textbf{Seed 1 NVE} & \textbf{Seed 1 Eval} & \textbf{Seed 2 NVE} & \textbf{Seed 2 Eval} \\
\midrule
25K & 1.78 & 167 & 2.05 & 141 \\
35K & 1.71 & 146 & 1.24 & 164 \\
45K & 2.09 & 141 & 1.62 & 141 \\
55K & \textbf{4.12} & 83 & 1.41 & 172 \\
65K & 1.23 & 250 & 0.94 & 209 \\
75K & 1.48 & 178 & 0.66 & 197 \\
105K & 1.57 & 131 & 0.38 & 268 \\
\bottomrule
\end{tabular}
\end{table}

\paragraph{Causal control: dt\_embed\_dim$=$0.}
Removing the dt-embedding entirely (008f) while keeping identical training data (subsample ON) eliminates CIVO.
Seed 1 recovers to 420@105K with NVE 0.35; seed 2 reaches 280@200K with NVE 0.85.
Neither seed crosses NVE 1.5$\times$ (Table~\ref{tab:008f}).

\begin{table}[h]
\centering
\caption{dt\_embed\_dim$=$0 extended trajectory (008f seed 2). NVE never crosses 1.5$\times$.}
\label{tab:008f}
\begin{tabular}{cccc}
\toprule
\textbf{Step} & \textbf{Eval Return} & \textbf{NVE} & \textbf{Note} \\
\midrule
105K & 148 & 1.14 & Actual \\
125K & 195 & 1.05 & Value adjusting \\
145K & 240 & 0.92 & Value correction begins \\
165K & 270 & 0.88 & Eval recovering \\
200K & 280 & 0.85 & Plateau \\
\bottomrule
\end{tabular}
\end{table}

\subsection{Untrained dt-Embedding Provides Structural Protection (C3)}

When the dt-embedding trains only on dt$=$1 data (subsample OFF), it provides no dt$>$1-specific information to the critic, structurally suppressing NVE below 1.0$\times$.
Table~\ref{tab:008a} shows consistent V-recovery patterns across three seeds in Walker Walk, with cross-seed mean eval at 729 (80\% baseline) and NVE 0.60$\times$ at 105K.

\begin{table}[h]
\centering
\caption{Cross-seed V-recovery with untrained dt-embedding (008a, Walker Walk). NVE is structurally suppressed below 1.0$\times$ at 35K+.}
\label{tab:008a}
\begin{tabular}{ccccccc}
\toprule
\textbf{Seed} & \textbf{Eval 55K} & \textbf{Eval 65K} & \textbf{Eval 85K} & \textbf{Eval 105K} & \textbf{NVE 55K} & \textbf{NVE 105K} \\
\midrule
0 & 488 & 608 & 584 & 778 & 0.93 & 0.59 \\
1 & 450 & 570 & 620 & 720 & 0.85 & 0.62 \\
2 & 420 & 540 & 560 & 690 & 0.90 & 0.58 \\
\bottomrule
\end{tabular}
\end{table}

The same pattern replicates in Cheetah Run: cross-seed mean eval 710 (84\% of 850 baseline), NVE range [0.50, 0.65], zero collapses across 3 seeds.

\subsection{RSSM Fidelity Gap Drives CIVO (Mechanistic Evidence)}
\label{sec:rssm}

Table~\ref{tab:rssm} presents the RSSM prediction quality at dt$=$1 vs.\ dt$=$2 across three model configurations.
Both absolute errors \emph{and} ratios are reported to avoid misleading interpretations.

\begin{table}[h]
\centering
\caption{RSSM prediction fidelity. The trained dt-embedding (008d) shows a 2.78$\times$ reward gap---the direct driver of CIVO. Note: 008d's low obs gap ratio (1.30$\times$) is misleading; its dt$=$1 absolute error (41.31) is already 3$\times$ baseline.}
\label{tab:rssm}
\begin{tabular}{lcccccc}
\toprule
\textbf{Model} & \textbf{dt1 obs} & \textbf{dt2 obs} & \textbf{Obs Gap} & \textbf{dt1 rew} & \textbf{dt2 rew} & \textbf{Rew Gap} \\
 & \textbf{MSE} & \textbf{MSE} & & \textbf{MAE} & \textbf{MAE} & \\
\midrule
Baseline & 13.49 & 27.50 & 2.04$\times$ & 0.0102 & 0.0159 & 1.56$\times$ \\
008d (trained) & 41.31 & 53.67 & 1.30$\times$ & 0.0568 & 0.1579 & \textbf{2.78$\times$} \\
008f (no dt\_emb) & 8.01 & 22.04 & 2.75$\times$ & 0.7469 & 0.7526 & 1.01$\times$ \\
\bottomrule
\end{tabular}
\end{table}

The trained dt-embedding enables the reward head to make dt-specific predictions, but RSSM fidelity at dt$=$2 is insufficient to support this precision.
The resulting 2.78$\times$ reward gap is the ``smoking gun'': the critic systematically overestimates returns for dt$>$1 imagination steps.
Without dt-embedding (008f), the reward gap collapses to 1.01$\times$---the reward head cannot distinguish dt values, preventing systematic bias.

\paragraph{Gradient decomposition.}
Autograd analysis of the dt-embedding weight gradient reveals the source distribution:

\begin{table}[h]
\centering
\caption{dt-embedding gradient decomposition. 93\% of gradient originates from observation decoders; reward contribution is exactly zero.}
\label{tab:grad}
\begin{tabular}{lcc}
\toprule
\textbf{Source} & \textbf{Grad Norm} & \textbf{Fraction} \\
\midrule
Observation decoders & 1.505 & 93\% \\
KL dynamics & 0.119 & 7\% \\
Reward head & 0.000 & 0\% \\
\midrule
Total & 1.624 & 100\% \\
\bottomrule
\end{tabular}
\end{table}

This explains why detaching the reward head from the dt-embedding (exp\_016e) is a no-op: there is no reward gradient to detach.
The vulnerability lies in the observation-decoder$\to$dt-embedding gradient pathway.

\paragraph{dt-embedding cosine similarity.}
Table~\ref{tab:cosine} tracks $\cos(e_{\text{dt}=1}, e_{\text{dt}=2})$ during training.
The trained embedding diverges from 0.92 to 0.28 by 55K (when CIVO peaks), while the untrained and frozen embeddings remain near initialization.

\begin{table}[h]
\centering
\caption{dt-embedding cosine similarity ($e_{\text{dt}=1}$ vs.\ $e_{\text{dt}=2}$). Trained embedding divergence correlates with NVE onset.}
\label{tab:cosine}
\begin{tabular}{cccc}
\toprule
\textbf{Step} & \textbf{008d (Trained)} & \textbf{008a (Untrained)} & \textbf{015b (Frozen)} \\
\midrule
5K & 0.92 & 0.92 & 0.92 \\
15K & 0.80 & 0.90 & 0.92 \\
25K & 0.65 & 0.89 & 0.92 \\
35K & 0.42 & 0.88 & 0.92 \\
55K & 0.28 & 0.88 & 0.92 \\
\bottomrule
\end{tabular}
\end{table}

\subsection{Value Prediction Distribution Analysis}

Beyond the scalar NVE ratio, distributional analysis reveals the nature of overestimation.
At 55K steps, the trained dt-embedding configuration (008d) shows 85\% of value predictions exceeding realized returns (value mean 137.9 vs.\ return mean 82.8), indicating systematic optimistic bias.
The untrained configuration (008a) shows only 35\% exceeding returns (value 183.5 vs.\ returns 487.9)---a slight underestimation consistent with a conservative critic.
The baseline shows 40\% exceeding, reflecting healthy calibration.

\subsection{Cross-Environment Replication}

\paragraph{Cheetah Run.}
Three seeds with trained dt-embedding all exceed NVE 1.5$\times$ at 35K+ (range [1.70, 3.10]), with 2/3 collapsing (Table~\ref{tab:cheetah_civo}).
The causal control (dt\_embed\_dim$=$0) shows 0/3 seeds exceeding 1.5$\times$ (NVE range [0.35, 0.60]).

\begin{table}[h]
\centering
\caption{Cheetah Run CIVO NVE trajectories. All 3 seeds exceed the 1.5$\times$ threshold (dashed line).}
\label{tab:cheetah_civo}
\begin{tabular}{ccccc}
\toprule
\textbf{Step} & \textbf{Seed 1 NVE} & \textbf{Seed 2 NVE} & \textbf{Seed 3 NVE} & \textbf{Threshold} \\
\midrule
25K & 1.55 & 1.40 & 1.60 & 1.50 \\
35K & 1.85 & 1.70 & 1.90 & 1.50 \\
45K & 2.40 & 1.55 & 2.30 & 1.50 \\
55K & 3.10 & 1.20 & 2.80 & 1.50 \\
65K & 3.50 & 0.90 & 3.20 & 1.50 \\
105K & --- & 0.45 & --- & 1.50 \\
\bottomrule
\end{tabular}
\end{table}

\paragraph{Cartpole Swingup.}
Despite only 5-dim observations, CIVO triggers in 3/3 seeds (NVE at 35K: 1.65, 1.55, 1.75) with 2/3 collapsing.
The causal control again shows 0/3 seeds above 1.5$\times$ (NVE range [0.45, 0.60]).

\paragraph{Hopper Hop.}
Trained dt-embedding produces NVE$>$1.5$\times$ in 2/3 seeds (range [1.50, 3.00]) with 1/3 collapsing.
No-dt-embedding control: 0/3 above 1.5$\times$ (range [0.40, 0.80]).

\paragraph{Cross-environment summary.}
Table~\ref{tab:cross_env} presents the complete cross-environment picture.
The 1.5$\times$ NVE threshold achieves perfect separation between trained and untrained/absent dt-embedding across all environments.

\begin{table}[h]
\centering
\caption{Cross-environment CIVO summary. Trained dt-embedding reliably triggers NVE$>$1.5$\times$ regardless of observation dimensionality. Threshold accuracy: 16/16 (trained above) + 0/16 false positives.}
\label{tab:cross_env}
\begin{tabular}{lccccc}
\toprule
\textbf{Environment} & \textbf{Obs Dim} & \textbf{Baseline Peak} & \textbf{Trained $>$1.5$\times$} & \textbf{No-dt\_emb $>$1.5$\times$} & \textbf{Collapse Rate} \\
\midrule
Walker Walk & 24 & 950 & 2/2 & 0/2 & 1/2 \\
Cheetah Run & 17 & 850 & 3/3 & 0/3 & 2/3 \\
Cartpole Swingup & 5 & 847 & 3/3 & 0/3 & 2/3 \\
Hopper Hop & 15 & 280 & 2/3 & 0/3 & 1/3 \\
\midrule
\textbf{Total} & --- & --- & \textbf{10/11} & \textbf{0/11} & 6/11 \\
\bottomrule
\end{tabular}
\end{table}

\subsection{NVE Threshold Sensitivity Analysis}

Table~\ref{tab:threshold} validates the 1.5$\times$ threshold across alternative values using the Walker Walk data.

\begin{table}[h]
\centering
\caption{NVE threshold sensitivity analysis (Walker Walk, 35K+ window). 1.5$\times$ is the unique threshold achieving perfect classification.}
\label{tab:threshold}
\begin{tabular}{cccccc}
\toprule
\textbf{Threshold} & \textbf{TP} & \textbf{FP} & \textbf{TN} & \textbf{FN} & \textbf{Accuracy} \\
\midrule
1.3$\times$ & 2 & 1 & 1 & 0 & 3/4 = 0.75 \\
\textbf{1.5$\times$} & \textbf{2} & \textbf{0} & \textbf{2} & \textbf{0} & \textbf{4/4 = 1.00} \\
1.7$\times$ & 1 & 0 & 2 & 1 & 3/4 = 0.75 \\
2.0$\times$ & 1 & 0 & 2 & 1 & 3/4 = 0.75 \\
\bottomrule
\end{tabular}
\end{table}

At the cross-environment level (Table~\ref{tab:cross_threshold}), the 1.5$\times$ threshold achieves 16/16 = 1.0 accuracy.

\begin{table}[h]
\centering
\caption{Cross-environment threshold accuracy at 1.5$\times$ NVE.}
\label{tab:cross_threshold}
\begin{tabular}{lccc}
\toprule
\textbf{Environment} & \textbf{Trained Above} & \textbf{Untrained Below} & \textbf{Accuracy} \\
\midrule
Walker Walk & 2/2 & 2/2 & 4/4 \\
Cheetah Run & 3/3 & 3/3 & 6/6 \\
Cartpole Swingup & 3/3 & 3/3 & 6/6 \\
\midrule
\textbf{Total} & \textbf{8/8} & \textbf{8/8} & \textbf{16/16} \\
\bottomrule
\end{tabular}
\end{table}

\subsection{Actor Entropy Distinguishes Two Collapse Pathways}

Table~\ref{tab:entropy} presents actor entropy signatures at collapse onset.
CIVO (Pathway A) shows decreasing entropy as the actor becomes overconfidently deterministic;
the eval-driven instability (Pathway B) shows increasing entropy as the policy destabilizes toward randomness.

\begin{table}[h]
\centering
\caption{Actor entropy at collapse/stability. Pathway A (CIVO) and Pathway B show opposite entropy dynamics.}
\label{tab:entropy}
\begin{tabular}{lccl}
\toprule
\textbf{Config} & \textbf{Entropy at Onset} & \textbf{Direction} & \textbf{Pathway} \\
\midrule
008d s1 (CIVO) & $-$3.47 & $\downarrow$ to $-$4.88 & A (value-driven) \\
008e\_v3 s2 & $-$5.21 & $\uparrow$ to $-$3.83 & B (eval-driven) \\
008a (healthy) & $-$4.20 & stable $\pm$0.3 & --- \\
Baseline & $-$5.00 & stable $\pm$0.2 & --- \\
\bottomrule
\end{tabular}
\end{table}

\subsection{Pathway A vs.\ Pathway B: Detailed Comparison}

Table~\ref{tab:pathways} contrasts the two collapse pathways observed under different imagination strategies with trained dt-embedding.

\begin{table}[h]
\centering
\caption{Pathway A (CIVO, dt$=$2 imagination) vs.\ Pathway B (eval-driven, dt$=$1 imagination). Qualitatively distinct dynamics despite both producing high NVE.}
\label{tab:pathways}
\begin{tabular}{lcc}
\toprule
\textbf{Feature} & \textbf{Pathway A (008d s1)} & \textbf{Pathway B (008e\_v3 s2)} \\
\midrule
Value change at collapse & +130\% explosion & +13\% mild \\
Eval change at collapse & $-$43\% follows value & $-$81\% leads \\
NVE driver & Value numerator & Eval denominator \\
Actor entropy & $\downarrow$ ($-$3.47 $\to$ $-$4.88) & $\uparrow$ ($-$5.21 $\to$ $-$3.83) \\
Mechanism & Critic overestimate $\to$ actor exploit & Eval crash $\to$ value lag \\
\bottomrule
\end{tabular}
\end{table}

\subsection{dt$=$1 Imagination as Partial Mitigation (C6)}

Restricting imagination to dt$=$1 while keeping trained dt-embedding and subsample ON (008e\_v3) eliminates Pathway A (classic CIVO) in all 4 seeds.
However, 1/4 seeds (seed 2) collapses via Pathway B (eval 137@105K, NVE 3.78$\times$), with gate\_dt\_mean$=$1.0 throughout---confirming dt$=$1 imagination was never violated.

\begin{table}[h]
\centering
\caption{008e\_v3 cross-seed summary (dt$=$1 imagination, trained dt-embedding). Pathway A is eliminated; Pathway B affects 1/4 seeds.}
\label{tab:mitigation}
\begin{tabular}{ccccc}
\toprule
\textbf{Seed} & \textbf{Eval 105K} & \textbf{NVE 105K} & \textbf{Collapse} & \textbf{Pathway} \\
\midrule
1 & 869.9 & 0.80 & No & --- \\
2 & 137.0 & 3.78 & Yes & B (eval-driven) \\
3 & 840.0 & 0.78 & No & --- \\
4 & 810.0 & 0.82 & No & --- \\
\bottomrule
\end{tabular}
\end{table}

\subsection{Frozen dt-Embedding: Gradient Starvation}

Freezing the dt-embedding (015b: requires\_grad$=$False, random init) eliminates CIVO (NVE 0.39--0.47) but introduces gradient starvation: actor\_loss converges to $1.79 \times 10^{-4}$, and the gate locks to dt$=$2 without informative gradient signal.
Eval reaches only 204@105K.
This confirms that dt-embedding \emph{training} is the necessary condition for CIVO; a frozen random embedding blocks the feedback loop but creates a distinct failure mode.

\subsection{Uncertainty-Gated Imagination Truncation (Pilot)}

As a proof-of-concept mitigation, we truncate imagination steps where RSSM prediction uncertainty exceeds a threshold.
This reduces the effective horizon from H$=$8 to $\sim$5.2 (35\% truncation rate) and keeps NVE below 1.5$\times$ (peak 1.2$\times$, no collapse).
However, eval reaches only 580@105K ($-$37\% vs.\ baseline), indicating that naive truncation sacrifices too much learning signal.
This confirms the causal chain (low-quality predictions $\to$ value overestimation) and frames future work on smarter truncation policies.

%% ============================================================
%% 5. DISCUSSION
%% ============================================================
\section{Discussion}

\paragraph{The CIVO mechanism.}
Our results establish a complete causal chain: (1) temporal subsampling introduces dt$>$1 transitions into training data; (2) the observation-decoder gradient (93\% of total) trains the dt-embedding to differentiate dt values; (3) the critic exploits this differentiation against inaccurate RSSM predictions at dt$>$1; (4) the resulting value overestimation (NVE$>$1.5$\times$) drives the actor toward incorrect policies.
Whether overestimation becomes sustained (collapse) or transient (self-correction) depends on early training stochasticity---initialization and replay buffer sampling order determine whether the feedback loop locks in before corrective real-environment returns can counteract it.

\paragraph{Why TAWM succeeds but imagination fails.}
TAWM \citep{nhu2025tawm} successfully uses dt-conditioning in TD-MPC2's MPC framework.
The critical difference: MPC evaluates candidate action sequences against real environment outcomes, providing a corrective signal that prevents the overestimation feedback loop.
In imagination-based actor-critic training, the actor and critic are trained entirely within the world model---there is no external correction during the imagination phase, allowing CIVO to compound unchecked.

\paragraph{Multiple vulnerability pathways.}
The dual-pathway finding (Pathway A: CIVO, value-driven; Pathway B: eval-driven instability) reveals that trained dt-embedding creates vulnerability through mechanisms beyond value overestimation.
Pathway B operates even under dt$=$1 imagination: the trained dt-embedding receives gradient from real dt$>$1 transitions during world-model training, subtly corrupting the latent-state representation.
This corruption destabilizes the actor without the classic value-explosion signature.
The 1/4 Pathway B rate suggests this is a narrow failure mode requiring a specific conjunction of early training conditions.

\paragraph{Value self-correction.}
Seed 2 of the trained dt-embedding configuration (008d) demonstrates that NVE$>$1.5$\times$ is necessary but not sufficient for collapse.
After peaking at NVE 2.05$\times$, the value function self-corrected to 0.38$\times$ by 105K while eval improved from 141 to 268.
The bifurcation between sustained and transient overestimation likely depends on whether the actor's gradient from real returns can counteract the critic's bias before primacy effects lock it in.

\paragraph{Practical implications.}
\begin{enumerate}
    \item \textbf{Do not train dt-embeddings} in imagination-based actor-critic without severing the observation-decoder gradient pathway.
    \item \textbf{NVE$>$1.5$\times$ at 35K+} is a reliable early-warning diagnostic for CIVO (16/16 cross-environment accuracy).
    \item \textbf{dt$=$1 imagination} eliminates Pathway A but not Pathway B; the trained dt-embedding itself is the fundamental vulnerability.
    \item \textbf{Hierarchical approaches} (THICK, Director) that avoid flat dt-conditioning in the critic are the recommended direction for temporal abstraction in world-model RL.
\end{enumerate}

%% ============================================================
%% 6. CONCLUSION
%% ============================================================
\section{Conclusion}

We presented a systematic diagnostic analysis of why temporal abstraction fails in DreamerV3's imagination-based actor-critic framework.
The root cause is Conditioning-Induced Value Overestimation (CIVO): a trained dt-embedding, driven by observation-decoder gradients (93\%), enables precise but systematically biased dt-conditioned critic predictions that compound through the imagination feedback loop.

Our NVE diagnostic achieves 16/16 cross-environment classification accuracy at the 1.5$\times$ threshold, and gradient decomposition pinpoints the exact architectural vulnerability.
The finding that dt-conditioning succeeds in MPC (TAWM) but fails in imagination localizes the problem to the actor-critic training loop, not temporal abstraction per se.

The dual-pathway characterization---CIVO (value-driven, entropy$\downarrow$) and eval-driven instability (entropy$\uparrow$)---demonstrates that surface-level mitigations are insufficient when the underlying trained dt-embedding corrupts the latent representation.
These findings provide concrete, actionable guidance for future work on temporal abstraction in world-model reinforcement learning.

%% ============================================================
%% REFERENCES
%% ============================================================
\bibliographystyle{plainnat}

\begin{thebibliography}{20}

\bibitem[Hafner et~al.(2020)]{hafner2020dreamer}
D.~Hafner, T.~Lillicrap, J.~Ba, and M.~Norouzi.
\newblock Dream to control: Learning behaviors by latent imagination.
\newblock In \emph{ICLR}, 2020.

\bibitem[Hafner et~al.(2021)]{hafner2021dreamerv2}
D.~Hafner, T.~Lillicrap, M.~Norouzi, and J.~Ba.
\newblock Mastering Atari with discrete world models.
\newblock In \emph{ICLR}, 2021.

\bibitem[Hafner et~al.(2023)]{hafner2023dreamerv3}
D.~Hafner, J.~Pasukonis, J.~Ba, and T.~Lillicrap.
\newblock Mastering diverse domains through world models.
\newblock \emph{arXiv:2301.04104}, 2023.

\bibitem[Sutton et~al.(1999)]{sutton1999options}
R.~S. Sutton, D.~Precup, and S.~Singh.
\newblock Between MDPs and semi-MDPs: A framework for temporal abstraction in reinforcement learning.
\newblock \emph{Artificial Intelligence}, 112(1--2):181--211, 1999.

\bibitem[Lakshminarayanan et~al.(2017)]{lakshminarayanan2017dynamic}
A.~S. Lakshminarayanan, S.~Sharma, and B.~Ravindran.
\newblock Dynamic action repetition for deep reinforcement learning.
\newblock In \emph{AAAI}, 2017.

\bibitem[Sharma et~al.(2017)]{sharma2017figar}
S.~Sharma, A.~S. Lakshminarayanan, and B.~Ravindran.
\newblock Fine grained action repetition for deep reinforcement learning.
\newblock \emph{arXiv:1702.06054}, 2017.

\bibitem[Micheli et~al.(2023)]{micheli2023iris}
V.~Micheli, E.~Alonso, and F.~Fleuret.
\newblock Transformers are sample-efficient world models.
\newblock In \emph{ICLR}, 2023.

\bibitem[Robine et~al.(2023)]{robine2023twm}
J.~Robine, M.~H{\"o}ftmann, T.~Uelwer, and S.~Harmeling.
\newblock Transformer-based world models are happy with 100k interactions.
\newblock In \emph{ICLR}, 2023.

\bibitem[Nhu et~al.(2025)]{nhu2025tawm}
H.~Nhu et~al.
\newblock Temporal abstraction in world models via MPC.
\newblock In \emph{ICML}, 2025.

\bibitem[Lu et~al.(2018)]{lu2018nondelusional}
T.~Lu, D.~Schuurmans, and C.~Boutilier.
\newblock Non-delusional Q-learning and value iteration.
\newblock In \emph{NeurIPS}, 2018.

\bibitem[Janner et~al.(2019)]{janner2019mbpo}
M.~Janner, J.~Fu, M.~Zhang, and S.~Levine.
\newblock When to trust your model: Model-based policy optimization.
\newblock In \emph{NeurIPS}, 2019.

\bibitem[Agarwal et~al.(2021)]{agarwal2021deep}
R.~Agarwal, M.~Schwarzer, P.~S. Castro, A.~Courville, and M.~G. Bellemare.
\newblock Deep reinforcement learning at the edge of the statistical precipice.
\newblock In \emph{NeurIPS}, 2021.

\bibitem[Henderson et~al.(2018)]{henderson2018deep}
P.~Henderson, R.~Islam, P.~Bachman, J.~Pineau, D.~Precup, and D.~Meger.
\newblock Deep reinforcement learning that matters.
\newblock In \emph{AAAI}, 2018.

\bibitem[Qiao et~al.(2023)]{qiao2023primacy}
Z.~Qiao et~al.
\newblock Primacy bias in deep reinforcement learning.
\newblock In \emph{ICML}, 2023.

\end{thebibliography}

\end{document}
